% TOP_PROBLEM: {"id": 3100089, "title": "Let G be a bicolored graph, and let H be the graph whose vertices are the valid pressing sequences of G and whose edges", "category_id": 2, "category": {"id": 2, "name": "combinatorics", "display_name": "Combinatorics", "description": "Counting problems, graph theory, discrete structures.", "slug": "combinatorics", "order_index": 2, "created_at": "2026-07-31T15:26:25.671Z"}, "status": "open"}
% FIRST_POSTED: null
% ENTRY_KIND: statement_only
% Imported from: batch11.tex; UnsolvedMath ID 3100089
\documentclass[11pt]{article}
\usepackage{amsmath,amssymb,mathtools}
\usepackage{fontspec,unicode-math}
\setmainfont{Latin Modern Roman}
\setmathfont{Latin Modern Math}
\usepackage{xurl,hyperref}
\newcommand{\sref}[2]{\hyperlink{source:#1:#2}{[S#2]}}
\newcommand{\cataloguescope}[1]{\paragraph{Mathematical statement.}#1}
\begin{document}
% UnsolvedMath ID 3100089; source catalogue code AMR-030-0089
\setcounter{section}{3100088}
\section{Let G be a bicolored graph, and let H be the graph whose vertices are the valid pressing sequences of G and whose edges}\label{3100089}
{\small\sffamily UnsolvedMath ID 3100089 · Combinatorics}
\subsection{Definitions and mathematical statement}

\paragraph{Shared notation.}$\mathbb N=\{1,2,\ldots\}$, and $\mathbb Z,\mathbb Q,\mathbb R,\mathbb C$ denote the integers, rationals, reals and complex numbers. Any other conventions are specified locally.


A bicoloured graph is a finite simple graph $G$ whose vertices are coloured black or white. At a currently black vertex $v$, a press complements the induced simple graph on the closed neighbourhood $N[v]=\{v\}\cup\{w:\{v,w\}\in E(G)\}$ and reverses every colour in $N[v]$; other edges and colours are unchanged. A successful pressing sequence is a finite word in $V(G)$ whose successive presses are permitted and whose final graph is edgeless and all white. These are the sequences in the named Pressing Game Conjecture. One-letter edits in that conjecture mean insertion or deletion of one vertex symbol, with every such operation counting as one edit.
Define $H(G)$ to have the successful sequences as vertices, joining distinct sequences whenever at most four insertions and deletions transform one into the other. Intermediate words in computing this edit distance need not themselves be successful.

\paragraph{Statement recorded in the supplied source.}
For every bicoloured graph $G$, can any two successful pressing sequences be joined by a finite path in $H(G)$? Thus every successive pair along the path is successful and has insertion--deletion distance at most four. Connectivity here means this pairwise path condition, including the vacuous case in which there are no successful sequences. \sref{3100089}{2}

\subsection{Short English statement}
Can any successful pressing sequence of a bicoloured graph be transformed into any other through successful sequences, using at most four insertions or deletions at each step?
\subsection{Sources}
\begin{description}
\item[\hypertarget{source:3100089:1}{S1}] ULAM.AI and the UnsolvedMath Contributors, \emph{UnsolvedMath: A Curated Collection of Open Mathematics Problems}, record 3100089 (AMR-030-0089). CC BY 4.0. \url{https://huggingface.co/datasets/ulamai/UnsolvedMath}.
\item[\hypertarget{source:3100089:2}{S2}] Joshua Cooper and Jeffrey Davis, \emph{Successful pressing sequences for a bicolored graph and binary matrices} (2015 revision), Section1, Definitions1--2, successful-sequence definition and Conjecture1, pages2--3. \url{https://arxiv.org/abs/1502.07450}.
\end{description}
\label{end:3100089}
\end{document}
